\chapter*{Preface to the Revised Edition}

\vspace{-1em}
This new edition of \emph{Lectures in Diffeology} is not merely a correction of typos or a polishing of the English prose,
although it is certainly that as well.
It is an opportunity to revisit the foundational philosophy of the theory in light of recent investigations.

In the original preface,
I emphasized the lineage of the \emph{Erlangen Program},
echoing the view of Jean-Marie Souriau that ``the group is the geometry.''
This perspective holds that a geometric object is fully characterized by its symmetries.
While this remains a powerful guiding principle,
the flexibility of diffeology has recently forced us to refine this view.
It has revealed that the ``absolute'' Kleinian program—where the abstract group determines the space—is insufficient when we step outside the rigid world of manifolds.

This insight comes from a counter-example I recently detailed,
known as the \emph{Boman Paradox}.%
\footnote{P. Iglesias-Zemmour, \emph{The Boman Paradox: When topology plus algebra do not define geometry}, 2025.}
Consider the \emph{Wire Plane} $\cW$,
a space whose underlying set is $\RR^2$ but whose plots are restricted to those that locally factor through smooth curves.
We have proven that this space shares the exact same topology as the standard Euclidean plane,
and remarkably,
its group of diffeomorphisms $\Diff(\cW)$ is algebraically identical to the standard group $\Diff(\RR^2)$.
Yet,
these spaces are geometrically distinct:
the standard plane is 2-dimensional,
while the Wire Plane is strictly 1-dimensional.

This phenomenon demonstrates that Geometry cannot be reduced to the sum of Topology and Algebra.
The local smooth structure—the definition of the \emph{plots}—precedes the global symmetry.
The space determines the group,
but the abstract group is not enough to reconstruct the space.
To resolve the paradox,
one must look at the group $\Diff(\cW)$ not as a set,
but as a \emph{diffeological group}.
Only then does the difference appear:
the isomorphism between the two groups is algebraic,
but it is not smooth.
Moreover,
the space itself is not homogeneous under the action of its own automorphisms.

Thus,
diffeology fulfills the Kleinian dream but with a crucial caveat:
the invariants are not carried by the group as an abstract algebraic entity,
but by the group as a diffeological space itself.
This reinforces the central thesis of these lectures:
that the \emph{diffeological structure} is the primary object,
capable of capturing subtleties—like the collapse of dimension in the Wire Plane—that coarser theories miss.

The paradox is developed in full detail in two new chapters,
B26 and B27,
where we prove that the D-topology of $\cW$ coincides with the Euclidean topology (Chapter B26)
and that $\Diff(\cW)$ and $\Diff(\RR^2)$ are distinct as diffeological groups (Chapter B27).

Beyond this epistemological update,
this revision includes several substantial additions and improvements:

\smuline{New chapters.}
Three entirely new chapters have been added to the Lectures.
The first,
Chapter L11 on \emph{Geometric Quantization by Paths},
describes the diffeological construction of the prequantum groupoid
(a natural extension of the classical Kostant-Souriau geometric quantization)
as a diffeological quotient of the space of paths of a parasymplectic space.
This shows how the quantum phase emerges naturally from this quotient
and demonstrates that the group of quantum automorphisms of the prequantum groupoid coincides exactly with the group of automorphisms of the parasymplectic structure,
without central extension at this level.

The other two new chapters,
B26 and B27,
develop the \emph{Boman Paradox} mentioned above,
providing the complete proofs of the topological identity and the group-theoretic distinction.

\smuline{Corrections and improvements.}
Beyond the addition of new material,
the entire book has been carefully revised.
Numerous typos,
misprints,
and grammatical inconsistencies have been corrected throughout all chapters.
Notational ambiguities have been resolved.
The preface itself has been updated to reflect the refined philosophical stance on the Kleinian--Souriau program
in light of the Boman Paradox,
and the literary piece ``A Wonderland of Smoothness'' has been added to capture the magical spirit of diffeology.

Taken together,
these additions and corrections transform the \emph{Lectures in Diffeology} from a set of seminar notes into a more complete and polished companion to the textbook,
while preserving the informal,
exploratory spirit that characterized the original.

\medskip
\centerline{\Large \bf A Wonderland of Smoothness}
\smallskip

Diffeology opens a door.
Behind it lies a landscape that classical differential geometry could only glimpse from afar — a kind of mathematical Wonderland,
or perhaps the Emerald City of Oz,
where familiar rules give way to stranger,
richer geometries.

Consider the \emph{irrational torus}.
It is the Cheshire Cat of this world:
it has no local coordinates,
no atlas,
no chart to pin it down.
Its topology is trivial
— the D-topology is coarse, almost asleep —
yet its geometry is dense,
intricate,
and alive.
The cat smiles,
but there is no body to hold the smile.
The torus exists,
but not as a manifold.

Then there is the \emph{Wire Plane}.
From a distance
— topologically speaking —
it looks just like Kansas.
The same points, the same open sets,
the same continuous maps. But step inside,
and the rules change.
The dimension collapses from 2 to 1.
You can move along a wire,
but you cannot slide transversally.
The symmetries are the same set as those of the Euclidean plane,
yet the smooth paths are fewer.
This is not Kansas anymore.
This is Oz.

The \emph{semi-cubic} is the looking-glass itself.
It is embedded in the plane — a perfect image of the real line
— yet it is not smoothly embedded.
The cusp at the origin is a singularity that the subset diffeology cannot see,
but the ambient diffeomorphisms that preserve the curve can feel it.
Like Alice stepping through the mirror,
the semi-cubic reveals that embedding is not enough;
smoothness requires a deeper relation between the space and its surroundings.

The \emph{square} is the land of Oz in miniature.
The group of diffeomorphisms sees what pure topology cannot:
corners,
edges,
interior
— three distinct orbits where a mere homeomorphism sees only boundary and interior.
The geometry is not in the points themselves but in how they are organized under symmetries.
This is Klein's program made manifest.

The \emph{cross}
— the union of the two axes —
is the Mad Hatter's tea party.
Differential forms on the cross are simply restrictions of ambient forms,
yet the cross is not a manifold.
The rules are simple,
almost trivial,
yet the object is singular.
A tea party where everyone is seated,
but the chairs are arranged in a way that defies expectation.

The \emph{space of geodesics of the torus} is Wonderland itself.
Over the circle of directions, the fibers are sometimes circles,
sometimes lines,
mixed together densely.
A closed 2-form descends to the quotient,
creating a parasymplectic structure where no classical symplectic geometry could survive.
The irrational geodesics are the White Rabbit:
always running,
never repeating,
dense in the space yet never settling.

And then there is the \emph{vague adjunction of a point}
— the rabbit hole.
Add a point that counts for nothing,
whose only neighborhood is the whole space.
A point that is closed, yet topologically inseparable from the rest.
A point that changes everything by adding nothing.

Diffeology is not merely a generalization of differential geometry.
It is a liberation.
It says: \emph{you are no longer bound by the tyranny of coordinates. You can explore spaces that were previously invisible. You can do geometry anywhere, on any set, no matter how singular.}

Welcome to Wonderland.
The tea party is about to begin.

\smallskip
\begin{flushright}
  P. I.-Z. \\
  Jerusalem, May 2026.
\end{flushright}
